% =====================================================================
%  Chương IV — KẾT QUẢ THỰC NGHIỆM
%  File này được thiết kế để \input vào file luận văn chính.
%  Gói cần khai báo ở preamble của main.tex (đã dùng ở các chương trước):
%     \usepackage{amsmath, amssymb}
%     \usepackage{booktabs}
%     \usepackage{float}          % cho [H]
%     \usepackage{hyperref}
%  Số liệu được trích trực tiếp từ các file kết quả trong repo:
%     runs_ate/eval_joint_triplet_BERT.csv   (Joint, backbone BERT)
%     runs_ate/eval_joint_triplet_T5.csv     (Joint, backbone T5)
%     runs_ate/eval_results_BERT.csv         (Pipeline, backbone BERT)
%     runs_ate/eval_results_T5.csv           (Pipeline, backbone T5)
%     runs_ate/results_ate_incorrect.csv     (lỗi ATE)
% =====================================================================
\providecommand{\cmark}{\ensuremath{\checkmark}}
\providecommand{\xmark}{\ensuremath{\times}}

\chapter{IV. KẾT QUẢ THỰC NGHIỆM}
\label{chap:ketqua}

Chương này trình bày bộ dữ liệu, thiết lập thực nghiệm, kết quả định
lượng và phân tích. Mục~\ref{sec:kq_data} mô tả dữ liệu; mục~\ref{sec:kq_setup}
trình bày siêu tham số, không gian cấu hình và độ đo; mục~\ref{sec:kq_results}
báo cáo kết quả của hai kiến trúc Joint/Pipeline trên hai backbone
BERT/T5; mục~\ref{sec:kq_analysis} phân tích, so sánh; và
mục~\ref{sec:kq_error} phân tích lỗi. Mọi số liệu được trích trực tiếp từ
các file kết quả trong repo.

% =====================================================================
\section{Bộ dữ liệu}
\label{sec:kq_data}

\subsection{Định dạng và phân chia}

Dữ liệu là tập đánh giá khách sạn tiếng Việt, định dạng \texttt{.apc} với
mỗi mẫu gồm 4 dòng: câu (chứa placeholder \texttt{\$T\$}), aspect term,
aspect category, và cảm xúc. Tập được chia cố định thành train/dev/test
(Bảng~\ref{tab:kq_data_stats}); tập test không tham gia huấn luyện.

\begin{table}[H]
\centering
\caption{Thống kê các tập dữ liệu}
\label{tab:kq_data_stats}
\begin{tabular}{@{}lrl@{}}
\toprule
\textbf{Tập} & \textbf{Số mẫu} & \textbf{Vai trò} \\
\midrule
Train & 2\,448 & Huấn luyện (kèm dữ liệu bổ trợ \texttt{supplement}) \\
Dev   & 304   & Early stopping và chọn checkpoint tốt nhất \\
Test  & 312   & Đánh giá chính thức \\
\bottomrule
\end{tabular}
\end{table}

Không gian nhãn gồm 6 nhóm khía cạnh (\texttt{AMENITY}, \texttt{BRANDING},
\texttt{EXPERIENCE}, \texttt{FACILITY}, \texttt{LOYALTY}, \texttt{SERVICE})
và 3 cực cảm xúc (\textit{positive} / \textit{negative} / \textit{neutral}),
tạo thành các tổ hợp nhãn ghép dùng cho độ đo bộ ba.

\subsection{Mất cân bằng nhãn và dữ liệu bổ trợ}

Phân bố nhãn trong tập huấn luyện \textbf{mất cân bằng nghiêm trọng}:
tổ hợp \texttt{SERVICE\_positive} chiếm tỉ lệ rất cao, trong khi
\texttt{BRANDING\_neutral} và \texttt{FACILITY\_neutral} cực hiếm. Để giảm
mất cân bằng, tập huấn luyện được bổ sung hai file TSV
(\texttt{supplement/negative.tsv}, \texttt{supplement/neutral.tsv}, chỉ
gồm \texttt{text}/\texttt{term}/\texttt{sentiment}). Các mẫu bổ trợ được
đánh dấu \texttt{is\_supplement=True}: chúng \emph{chỉ} tham gia huấn
luyện đầu phân loại cảm xúc; đầu phân loại nhóm khía cạnh được che khỏi
chúng (mục~\ref{subsec:multitask_loss}, Chương~3). Như sẽ thấy ở
mục~\ref{sec:kq_error}, mức độ mất cân bằng này chính là nguyên nhân khiến
một số nhãn hiếm đạt F1 bằng 0 trên mọi cấu hình.

% =====================================================================
\section{Thiết lập thực nghiệm}
\label{sec:kq_setup}

\subsection{Siêu tham số}

\begin{table}[H]
\centering
\caption{Siêu tham số huấn luyện module APC}
\label{tab:kq_hparams}
\begin{tabular}{@{}lll@{}}
\toprule
\textbf{Siêu tham số} & \textbf{Giá trị} & \textbf{Ghi chú} \\
\midrule
Backbone        & \texttt{bert-base-uncased} / \texttt{t5-base} & encoder cho APC \\
Số epoch tối đa & 15  & \texttt{NUM\_EPOCHS} \\
Early stopping  & patience $= 4$ & theo dev joint F1 \\
Batch size      & 16  & \texttt{BATCH\_SIZE} \\
Learning rate   & $2\times10^{-5}$ & AdamW \\
Độ dài chuỗi    & 128 & \texttt{MAX\_SEQ\_LEN} \\
Dropout         & 0.1 & \texttt{DROPOUT} \\
Số đầu attention & 8  & \texttt{NUM\_HEADS} \\
SRD threshold $\alpha$ & 5 & bán kính ngữ cảnh LCF \\
ToMe merge steps & 2 & số vòng gộp (post-BERT) \\
Mixed precision  & Bật (AMP) & tự tắt nếu không có CUDA \\
Random seed      & 42 & cố định để tái lập \\
\bottomrule
\end{tabular}
\end{table}

Tham số được tối ưu bằng \textbf{AdamW}; sau mỗi epoch, mô hình được đánh
giá trên \texttt{dev} và lưu checkpoint có \textbf{joint F1 trên dev cao
nhất}, dừng sớm nếu không cải thiện sau \texttt{PATIENCE} epoch. Toàn bộ
thực nghiệm dùng chế độ \texttt{resize} (xem
mục~\ref{subsubsec:why_resize}, Chương~3).

\subsection{Không gian cấu hình}

Các cấu hình được tổ chức theo thiết kế nhân tố: bật/tắt LCF; biến thể LCF
(CDM/CDW); và chiến lược ToMe (\texttt{bipartite}/\texttt{sequential\_local}/%
\texttt{attention\_weighted}). Bảng~\ref{tab:kq_configs} giải thích quy ước
đặt tên 12 cấu hình.

\begin{table}[H]
\centering
\caption{Quy ước đặt tên cấu hình}
\label{tab:kq_configs}
\begin{tabular}{@{}lll@{}}
\toprule
\textbf{Nhóm cấu hình} & \textbf{LCF} & \textbf{ToMe} \\
\midrule
\texttt{baseline\_balanced}                 & \xmark & \xmark \\
\texttt{lcf\_only\_\{cdm,cdw\}}             & \cmark & \xmark \\
\texttt{\{bip,seq,attn\}\_resize}           & \xmark & \cmark (3 chiến lược) \\
\texttt{lcf\_\{bip,seq,attn\}\_\{cdm,cdw\}\_resize} & \cmark & \cmark (3 chiến lược $\times$ CDM/CDW) \\
\bottomrule
\end{tabular}
\end{table}

\subsection{Độ đo và quy trình đánh giá}

Các độ đo gồm: \textbf{ATE F1} (khớp chính xác); \textbf{Micro F1} và
\textbf{Macro F1} trên bộ ba (nhãn ghép \texttt{"<sentiment>\_<category>"},
TP đòi hỏi cả ba thành phần khớp gold); và các chỉ số \textbf{Oracle}
(Cat/Sent/Joint Acc) -- độ chính xác khi cho trước aspect term gold, phản
ánh giới hạn trên của bộ phân loại. Hai quy trình đánh giá:

\begin{itemize}
  \item \textbf{Joint}: ATE và APC đánh giá trên cùng đơn vị khớp với
        nhãn gold (ATE F1 $\approx 79.87$).
  \item \textbf{Pipeline}: T5 sinh aspect term tự do, APC phân loại trên
        các cụm từ \emph{được dự đoán} (ATE F1 $\approx 60$); phản ánh
        hiệu suất đầu-cuối thực tế kèm lan truyền lỗi.
\end{itemize}

% =====================================================================
\section{Kết quả thực nghiệm}
\label{sec:kq_results}

\subsection{Kiến trúc Joint}

Bảng~\ref{tab:kq_joint_bert} và~\ref{tab:kq_joint_t5} báo cáo 12 cấu hình
cho hai backbone (ATE F1 $= 79.87$ cho mọi cấu hình; cột thời gian là thời
gian huấn luyện, giây).

\begin{table}[H]
\centering
\caption{Kết quả Joint -- backbone \texttt{BERT} (ATE F1 $=79.87$)}
\label{tab:kq_joint_bert}
\small
\begin{tabular}{@{}lcccr@{}}
\toprule
\textbf{Cấu hình} & \textbf{Micro F1} & \textbf{Macro F1} & \textbf{Oracle Joint} & \textbf{Train (s)} \\
\midrule
baseline\_balanced       & 70.53 & 48.80 & 86.54 & 257 \\
lcf\_only\_cdw           & 69.24 & 46.66 & 85.90 & 265 \\
lcf\_only\_cdm           & 68.92 & 47.94 & 84.62 & 311 \\
bip\_resize              & 70.85 & 49.19 & 86.86 & 374 \\
lcf\_bip\_cdm\_resize    & 69.24 & 54.20 & 84.94 & 383 \\
lcf\_bip\_cdw\_resize    & 70.53 & 49.61 & 86.86 & 463 \\
lcf\_attn\_cdw\_resize   & 69.57 & 48.11 & 84.94 & 533 \\
\textbf{lcf\_seq\_cdm\_resize} & \textbf{71.18} & 50.00 & 86.86 & 553 \\
lcf\_attn\_cdm\_resize   & 69.89 & \textbf{57.16} & 86.54 & 586 \\
lcf\_seq\_cdw\_resize    & 70.21 & 50.00 & 86.86 & 608 \\
seq\_resize              & 69.24 & 47.98 & 86.22 & 651 \\
attn\_resize             & 68.28 & 47.76 & 85.26 & 684 \\
\bottomrule
\end{tabular}
\end{table}

\begin{table}[H]
\centering
\caption{Kết quả Joint -- backbone \texttt{T5} (ATE F1 $=79.87$)}
\label{tab:kq_joint_t5}
\small
\begin{tabular}{@{}lcccr@{}}
\toprule
\textbf{Cấu hình} & \textbf{Micro F1} & \textbf{Macro F1} & \textbf{Oracle Joint} & \textbf{Train (s)} \\
\midrule
baseline\_balanced       & 69.57 & 47.88 & 85.90 & 337 \\
lcf\_only\_cdw           & 69.89 & 48.76 & 86.54 & 345 \\
lcf\_only\_cdm           & 70.21 & 48.75 & 86.22 & 468 \\
bip\_resize              & 70.53 & 48.03 & 86.86 & 498 \\
lcf\_bip\_cdw\_resize    & 70.53 & 49.80 & 87.18 & 651 \\
attn\_resize             & 68.28 & 45.48 & 84.62 & 659 \\
lcf\_bip\_cdm\_resize    & 71.18 & \textbf{56.03} & 86.54 & 698 \\
lcf\_attn\_cdw\_resize   & 70.85 & 49.63 & 85.58 & 853 \\
\textbf{seq\_resize}     & \textbf{71.50} & 50.12 & \textbf{89.10} & 886 \\
lcf\_seq\_cdw\_resize    & 69.89 & 54.25 & 85.90 & 901 \\
lcf\_attn\_cdm\_resize   & 67.95 & 53.40 & 84.29 & 913 \\
lcf\_seq\_cdm\_resize    & 69.89 & 45.79 & 85.90 & 921 \\
\bottomrule
\end{tabular}
\end{table}

\subsection{Kiến trúc Pipeline}

Khi T5 sinh aspect term tự do (ATE F1 giảm còn $\approx 60$), Micro F1 bộ
ba giảm mạnh (Bảng~\ref{tab:kq_pipeline}). Cột \textit{Infer} là thời gian
suy luận APC trên các cụm từ dự đoán (giây).

\begin{table}[H]
\centering
\caption{Kết quả Pipeline -- một số cấu hình tiêu biểu (ATE F1: BERT $=59.76$, T5 $=60.00$)}
\label{tab:kq_pipeline}
\small
\begin{tabular}{@{}llccc r@{}}
\toprule
\textbf{Backbone} & \textbf{Cấu hình} & \textbf{Micro F1} & \textbf{Macro F1} & \textbf{Oracle Joint} & \textbf{Infer (s)} \\
\midrule
BERT & baseline\_balanced     & 52.86 & 38.41 & 86.54 & 10.8 \\
BERT & lcf\_seq\_cdm\_resize  & \textbf{53.57} & 40.29 & 86.86 & 14.9 \\
BERT & lcf\_attn\_cdm\_resize & 52.38 & \textbf{43.90} & 86.54 & 14.8 \\
\midrule
T5   & baseline\_balanced     & 52.86 & 39.57 & 85.90 & 7.6 \\
T5   & seq\_resize            & \textbf{53.81} & 39.98 & \textbf{89.10} & 14.3 \\
T5   & lcf\_bip\_cdm\_resize  & 53.10 & \textbf{47.81} & 86.54 & 11.3 \\
\bottomrule
\end{tabular}
\end{table}

% =====================================================================
\section{Phân tích và so sánh}
\label{sec:kq_analysis}

\subsection{Joint vượt Pipeline khoảng 18 điểm Micro F1}

Micro F1 bộ ba của Joint ($\approx 71$) cao hơn Pipeline ($\approx 53$)
khoảng 18 điểm. Nguyên nhân nằm hoàn toàn ở bước ATE: ATE F1 giảm từ
$79.87$ (Joint) xuống $\approx 60$ (Pipeline). Bằng chứng then chốt là
chỉ số \textbf{Oracle Joint Acc gần như không đổi} giữa hai kiến trúc với
cùng cấu hình (ví dụ \texttt{lcf\_attn\_cdm\_resize}: $86.54$ ở cả Joint
lẫn Pipeline) -- tức bộ phân loại là \emph{như nhau}, hạn chế đến từ chất
lượng aspect term được trích xuất, không phải từ classifier. Đây là biểu
hiện điển hình của \textbf{lan truyền lỗi} trong kiến trúc pipeline.

\subsection{BERT và T5 tương đương về Micro F1}

Hai backbone cho Micro F1 Joint sát nhau ($\approx 71$): tốt nhất là T5
\texttt{seq\_resize} ($71.50$) và BERT \texttt{lcf\_seq\_cdm\_resize}
($71.18$). T5 nhỉnh hơn về \textbf{Oracle Joint} ($89.10$ so với $86.86$),
cho thấy tiềm năng cao hơn \emph{nếu} ATE được cải thiện. Tuy nhiên, BERT
huấn luyện nhanh hơn rõ rệt (baseline $257$\,s so với $337$\,s; các cấu
hình LCF của T5 thường vượt $850$--$920$\,s).

\subsection{CDM tốt hơn CDW về Macro F1}

Trên nhãn thiểu số, mặt nạ nhị phân CDM bảo toàn ``tập trung cứng'' tốt
hơn trọng số mềm CDW, dẫn tới Macro F1 cao hơn đáng kể:

\begin{itemize}
  \item BERT, chiến lược \texttt{attention\_weighted}: CDM đạt Macro
        $57.16$ so với CDW $48.11$ ($+9.05$ điểm);
  \item T5, chiến lược \texttt{bipartite}: CDM đạt Macro $56.03$ so với
        CDW $49.80$ ($+6.23$ điểm).
\end{itemize}

Trong khi đó, Micro F1 giữa CDM và CDW chênh lệch nhỏ -- nghĩa là CDM cải
thiện chủ yếu ở các nhãn hiếm chứ không hy sinh nhãn phổ biến.

\subsection{Gộp token không làm giảm chất lượng}

Một kết quả đáng chú ý: các cấu hình có ToMe \textbf{không kém} baseline,
thậm chí tốt hơn nhẹ. Với BERT, baseline đạt Micro $70.53$ còn
\texttt{lcf\_seq\_cdm\_resize} đạt $71.18$ và \texttt{bip\_resize} đạt
$70.85$. Điều này khẳng định: việc gộp các token tương đồng (loại dư
thừa) \emph{không} phá hỏng biểu diễn cần cho phân loại -- tiền đề quan
trọng để áp dụng gộp token cho mục tiêu tăng tốc về sau.

\subsection{Chi phí huấn luyện trong chế độ resize}

Cần nhấn mạnh: trong chế độ \texttt{resize}, các cấu hình ToMe có thời
gian huấn luyện \textbf{cao hơn} baseline (ví dụ BERT: $257$\,s so với
$553$\,s ở \texttt{lcf\_seq\_cdm\_resize}), do thêm bước gộp và nội suy
\emph{mà không} rút ngắn chuỗi cho các lớp phía sau. Điều này một lần nữa
xác nhận: \textbf{resize chỉ phục vụ đo chất lượng biểu diễn, không phải
để tăng tốc}; lợi ích tốc độ thực sự thuộc về hướng phát triển
(Chương~5).

% =====================================================================
\section{Phân tích lỗi}
\label{sec:kq_error}

\subsection{Lỗi ở bước trích xuất khía cạnh (ATE)}

Khảo sát \texttt{results\_ate\_incorrect.csv} cho thấy lỗi ATE tập trung
ở hai dạng (Bảng~\ref{tab:kq_ate_err}): (i) \textbf{sai biên cụm từ} (trích
thừa/thiếu so với gold), và (ii) \textbf{chọn nhầm khía cạnh} trong câu
nhiều khía cạnh. Vì tiêu chí chấm là khớp chính xác, các lỗi biên này đều
bị tính sai hoàn toàn, góp phần kéo ATE F1 Pipeline xuống $\approx 60$.

\begin{table}[H]
\centering
\caption{Một số lỗi ATE tiêu biểu}
\label{tab:kq_ate_err}
\small
\begin{tabular}{@{}p{6.6cm}p{3.0cm}p{3.0cm}@{}}
\toprule
\textbf{Câu (rút gọn)} & \textbf{Dự đoán} & \textbf{Gold} \\
\midrule
\textit{Exceptional staff and customer service} & staff & staff and customer service \\
\textit{The service staff is very attentive and conscientious} & service staff & conscientious \\
\textit{they upgraded us to a beautiful suite\dots} & suite & they \\
\bottomrule
\end{tabular}
\end{table}

\subsection{Lỗi ở bước phân loại: nhãn thiểu số}

Ở mọi cấu hình, hai tổ hợp \texttt{BRANDING\_neutral} và
\texttt{FACILITY\_neutral} đạt \textbf{F1 $= 0$}, và \texttt{SERVICE\_negative}
thường rất thấp ($0$--$50$). Đây là hệ quả trực tiếp của mất cân bằng dữ
liệu (mục~\ref{sec:kq_data}): số mẫu của các tổ hợp này quá ít để mô hình
học. Chính các nhãn ``chết'' này kéo Macro F1 ($\approx 48$--$57$) xuống
thấp hơn nhiều so với Micro F1 ($\approx 71$). Đây là \textbf{vấn đề dữ
liệu}, không phải lỗi kiến trúc -- weighted loss và dữ liệu bổ trợ chỉ
giảm nhẹ chứ không khắc phục triệt để.

% =====================================================================
\section*{Tóm tắt chương}
\addcontentsline{toc}{section}{Tóm tắt chương}

Trên bộ dữ liệu đánh giá khách sạn tiếng Việt, kiến trúc \textbf{Joint}
đạt Micro F1 bộ ba tốt nhất $71.50$ (T5, \texttt{seq\_resize}) và $71.18$
(BERT, \texttt{lcf\_seq\_cdm\_resize}), vượt \textbf{Pipeline}
($\approx 53$) khoảng 18 điểm do lan truyền lỗi từ ATE -- được xác nhận
bởi Oracle Joint Acc gần như không đổi giữa hai kiến trúc. \textbf{CDM}
cải thiện Macro F1 trên nhãn thiểu số ($+6$ đến $+9$ điểm so với CDW), còn
\textbf{gộp token} không làm giảm (thậm chí cải thiện nhẹ) chất lượng so
với baseline. Phân tích lỗi chỉ ra hai nút thắt: lỗi biên/chọn nhầm ở ATE
và các nhãn thiểu số ``chết'' do mất cân bằng dữ liệu. Đáng chú ý, trong
chế độ \texttt{resize}, gộp token làm tăng (không giảm) thời gian huấn
luyện -- khẳng định resize chỉ đo chất lượng biểu diễn; tăng tốc thực sự
được bàn ở Chương~5.